    \subsection{Attribution graph}
    The attribution graph is a directed acyclic graph $G = (V, E, W)$ with
    vertex set partitioned into
    $V = V_{\text{emb}} \cup V_{\text{inter}} \cup V_{\text{error}} \cup V_{\text{logit}}$
    (embedding, intermediate, error, and logit nodes) and signed weights
    $W \in \mathbb{R}^{N \times N}$ on $N := |V|$ nodes, with $W_{ij}$ the
    linear influence of source node $j$ on target node $i$~\citep{ameisen2024scaling}.
    Embedding and error nodes are sources (no incoming edges), logit nodes
    are sinks (no outgoing edges), and each intermediate node carries signal
    between them, indexed by a transformer (layer, position) pair. A node has edges to every other nodes of higher layer.
    The reader also supplies a token weight $\tau \in \mathbb{R}^N$ with
    $\sum_e \tau_e = 1$ marking the input tokens of interest, and a logit
    weight $\ell \in \mathbb{R}^N$ marking the output positions of interest;
    both enter the loss terms below.


    \subsection{Summarization graph}
    A \emph{summary} of $G$ is a triple $(V', E', \varphi)$ comprising
    \begin{itemize}[leftmargin=*,topsep=2pt,itemsep=1pt]
    \item a pruned subgraph $V' \subseteq V$, $E' \subseteq E \cap (V' \times V')$;
    \item a partition $\varphi : V' \to \pi(\varphi)$ where
    $\pi(\varphi) := \{\varphi(v) : v \in V'\}$ is the set of nonempty
    \emph{supernodes} and $K := |\pi(\varphi)|$ their count. Every
    $v \in (V_{\text{emb}} \cup V_{\text{error}} \cup V_{\text{logit}}) \cap V'$ is a singleton
    supernode; only the feature nodes
    $V'_{\text{mid}} := V' \cap V_{\text{inter}}$ are free to be grouped.
    \end{itemize}
    The summary induces the \emph{summarization graph}
    $G_{SN} = (\pi(\varphi),\, E^{SN},\, W^{SN})$. For each ordered pair of
    distinct supernodes $(S, T)$, define the signed block sum
    \begin{equation}
    B_{S \to T}
    \;=\;
    \sum_{u \in S,\, v \in T} W_{vu}\,\mathbf{1}[(u, v) \in E']
    \label{eq:block-sum}
    \end{equation}
    (the signed weight of all edges from members of $S$ to members of $T$;
    $W_{vu}$ is the entry of $W$ for the edge $u \to v$, with target $v$
    indexing the row and source $u$ the column). The supernode
    adjacency keeps only the dominant signed direction per unordered pair:
    \begin{equation}
    W^{SN}_{S \to T}
    \;=\;
    \begin{cases}
    B_{S \to T} & \text{if } |B_{S \to T}| > |B_{T \to S}|, \\
    0           & \text{otherwise,}
    \end{cases}
    \label{eq:wsn}
    \end{equation}
    so $W^{SN}_{S \to T} \ne 0$ implies $W^{SN}_{T \to S} = 0$, eliminating
    $2$-cycles in $G_{SN}$ by construction. Longer cycles are broken by
    the projection step in Section~\ref{sec:solution-dag}, and we set
    $E^{SN} = \{(S, T) : W^{SN}_{S \to T} \ne 0\}$.


    \subsection{Flow quantities and role vectors}\label{sec:flow}
    Two propagated scores quantify how strongly each node participates in
    the $V_{\text{emb}} \to V_{\text{logit}}$ pathway. With the
    column-normalized weight matrix
    $W^{\text{in}}_{ij} = |W_{ij}|/\sum_k |W_{kj}|$ and the row-normalized
    $W^{\text{out}}_{ij} = |W_{ij}|/\sum_k |W_{ik}|$, the per-node
    \emph{influence} (downstream reach to logits) and \emph{relevance}
    (upstream reach to input tokens) are
    \begin{align}
    \mathrm{Inf} &= \ell^\top \bigl(I - (W^{\text{in}})^\top\bigr)^{-1}
    \;\in\; \mathbb{R}^N,
    \label{eq:inf}\\
    \mathrm{Rel} &= \tau^\top \bigl(I - W^{\text{out}}\bigr)^{-1}
    \;\in\; \mathbb{R}^N.
    \label{eq:rel}
    \end{align}
    The Neumann series converges because $G$ is a DAG. Edge-level
    analogues weight a directed edge $u \to v$ (matrix entry $W_{vu}$)
    by the importance of its endpoint:
    \begin{equation}
    \mathrm{Inf}_{u \to v}
    \;=\;
    \mathrm{Inf}_v \,\frac{|W_{vu}|}{\sum_k |W_{ku}|},
    \qquad
    \mathrm{Rel}_{u \to v}
    \;=\;
    \mathrm{Rel}_u \,\frac{|W_{vu}|}{\sum_k |W_{vk}|}.
    \label{eq:inf-rel-edge}
    \end{equation}
    The \emph{role vector} of a feature $u \in V'_{\text{mid}}$ stacks its
    weighted outgoing and incoming edge profiles, with each coordinate
    scaled by the importance of the corresponding endpoint:
    \begin{align}
    (v_u^{\text{out}})_k &= W_{ku}\,\sqrt{\mathrm{Inf}_k},
    \quad k \in V', \nonumber\\
    (v_u^{\text{in}})_k  &= W_{uk}\,\sqrt{\mathrm{Rel}_k},
    \quad k \in V', \nonumber\\
    r(u) &= \bigl[\,v_u^{\text{out}}\,;\, v_u^{\text{in}}\,\bigr]
    \;\in\; \mathbb{R}^{2|V'|}.
    \label{eq:role}
    \end{align}
    Outgoing coordinates are weighted by the downstream importance of the
    target ($\sqrt{\mathrm{Inf}_k}$), incoming coordinates by the upstream
    importance of the source ($\sqrt{\mathrm{Rel}_k}$), so $r(u)$
    concentrates on the edge dimensions that actually carry causal mass.


    \subsection{Desiderata}
    A reader inspects $G_{SN}$ to recover the model's causal story: the
    summary should be \emph{interpretable} (they can trace the flow of
    information at a glance) and \emph{faithful} (perturbing a supernode
    in $G_{SN}$ predicts the effect on other supernodes and on the
    logits). Three quality criteria translate these requirements into
    measurable conditions on the summary. Well-formedness of $\varphi$
    and acyclicity of $G_{SN}$ are not desiderata but properties of the
    construction above: the partition definition forces embedding,
    error, and logit nodes to be singletons, the dominant-direction
    rule of Eq.~\eqref{eq:wsn} eliminates $2$-cycles, and longer cycles
    are broken in Section~\ref{sec:solution-dag}.
    \begin{description}[leftmargin=2.0em,style=nextline,topsep=2pt,itemsep=2pt]
    \item[(D1) Complexity.] $G_{SN}$ has few enough supernodes to
    inspect in one sitting.
    \item[(D2) Atomicity.] Members of a supernode share a single causal
    role---similar influential downstream targets and similar relevant
    upstream sources---so the supernode reads as one mechanism, and the
    partition cannot be improved by further splitting or merging.
    \item[(D3) Causal information preservation.] $G_{SN}$ retains the
    edge weight that carries causal effect in the pruned graph. Edges
    absorbed inside a supernode and edge mass dropped by the
    dominant-direction rule both become invisible in $G_{SN}$: a reader
    using $G_{SN}$ to predict the outcome of a steering intervention on
    the underlying model has no edge to attribute the lost contribution
    to, even though the model still produces it. The partition should
    keep as much pruned-graph edge weight visible in $G_{SN}$ as
    possible.
    \end{description}

    \subsection{Objective}\label{sec:objective}
    We turn (D1)--(D3) into a constrained optimisation on the feasible
    set already pinned down by the summary-graph definition. Three soft
    losses---each in $[0, 1]$, one per desideratum---combine as a
    weighted sum on the probability simplex:
    \begin{equation}
    \min_{(V', E', \varphi)}\;
    L \;=\;
    \lambda_{\text{cplx}}\, L_{\text{cplx}}
    \;+\;
    \lambda_{\text{atom}}\, L_{\text{atom}}
    \;+\;
    \lambda_{\text{causal}}\, L_{\text{causal}},
    \quad
    \lambda_i \ge 0,\;\; \sum_i \lambda_i = 1.
    \label{eq:Lsum}
    \end{equation}
    We treat $\lambda_{\text{causal}} = 1 - \lambda_{\text{cplx}} -
    \lambda_{\text{atom}}$ as derived and report the remaining two.
    Each $L_i \in [0, 1]$, so $L \in [0, 1]$.

    \paragraph{Complexity loss (D1).}
    The total number of supernodes the reader must hold, normalised by
    the pruned graph size:
    \begin{equation}
    L_{\text{cplx}}
    \;=\;
    \frac{|\pi(\varphi)|}{|V'|}
    \;\in\; [0, 1].
    \label{eq:Lcplx}
    \end{equation}
    $L_{\text{cplx}}$ shrinks as feature nodes are merged and reaches
    $1$ when every node is its own singleton; the embedding, error, and
    logit nodes are always singletons by construction.

    \paragraph{Atomicity loss (D2).}
    With role distance $d(u, v) = 1 - \cos_{+}\!\bigl(r(u), r(v)\bigr)$ on
    the role vectors of Eq.~\eqref{eq:role}, where
    $\cos_{+}(x, y) = \max(0, \cos(x, y))$, the silhouette of a feature
    $u \in V'_{\text{mid}}$ is
    \begin{equation}
    s(u)
    \;=\;
    \frac{b(u) - a(u)}{\max\bigl(a(u),\, b(u)\bigr)}
    \;\in\; [-1, 1],
    \label{eq:silhouette}
    \end{equation}
    with
    \[
    a(u) = \tfrac{1}{|\varphi(u)| - 1}
    \sum_{u' \in \varphi(u),\, u' \ne u}\!\! d(u, u'),
    \qquad
    b(u) = \min_{S \ne \varphi(u)}\, \tfrac{1}{|S|}\!\sum_{u' \in S}\! d(u, u'),
    \]
    and the convention $s(u) = -1$ when $|\varphi(u)| = 1$, so unconsolidated
    singletons are penalised maximally. The atomicity loss is the
    mean-silhouette deficit,
    \begin{equation}
    L_{\text{atom}}
    \;=\;
    \frac{1 - \bar{s}}{2},
    \qquad
    \bar{s} \;=\; \frac{1}{|V'_{\text{mid}}|}
    \sum_{u \in V'_{\text{mid}}} s(u),
    \label{eq:Latom}
    \end{equation}
    in $[0, 1]$. The silhouette form encodes ``the partition cannot be
    improved by splitting or merging'' directly: $s(u) > 0$ iff $u$ is
    closer to its own supernode than to the nearest other, so neither
    moving $u$ nor splitting $\varphi(u)$ would locally improve the
    assignment. $L_{\text{atom}}$ rises when members of a supernode are
    dissimilar (over-merging) \emph{and} when similar features sit in
    different supernodes (under-merging).

    \paragraph{Causal-preservation loss (D3).}
    Three mechanisms remove pruned-graph edge weight from $G_{SN}$:
    (i) an edge $u \to v$ with $u, v$ in the same supernode disappears
    entirely; (ii) on each unordered pair $\{S, T\}$ with flow in both
    directions, the dominant-direction rule of Eq.~\eqref{eq:wsn}
    discards the smaller of $|B_{S \to T}|, |B_{T \to S}|$;
    (iii) opposing-sign edges from $S$ to $T$ partially cancel inside
    $|B_{S \to T}|$. All three hide pathway weight from $G_{SN}$: a
    reader using $G_{SN}$ to predict the result of a steering
    intervention has no edge to attribute the lost contribution to,
    though the underlying model still produces it. The loss is the
    fraction of pruned edge weight not visible in $G_{SN}$:
    \begin{equation}
    L_{\text{causal}}
    \;=\;
    1 \;-\;
    \frac{\sum_{S \ne T}\, |W^{SN}_{S \to T}|}
         {\sum_{(u, v) \in E'} |W_{vu}|}
    \;\in\; [0, 1].
    \label{eq:Lcausal}
    \end{equation}
    $L_{\text{causal}} = 0$ when every pruned-graph edge survives intact
    (no intra-supernode edges, no $2$-cycles between supernodes, no
    sign cancellation within any block sum); it grows when supernodes
    swallow direct edges or span overlapping layer ranges. Because
    $L_{\text{atom}}$ rewards merging role-similar features even when
    they are connected by a strong direct edge, $\lambda_{\text{causal}}$
    provides the explicit knob for the \emph{diamond} case where keeping
    that direct edge visible matters more than collapsing the pair.

    \paragraph{Pruning is scored separately.}
    Stage~1 chooses $V'$ to retain the causal mass flowing from
    $V_{\text{emb}}$ to $V_{\text{logit}}$ in $G$. With $\mathrm{Rel}$ of
    Eq.~\eqref{eq:rel} recomputed on the pruned graph (denoted
    $\mathrm{Rel}_{V'}$), the \emph{flow completeness}
    \begin{equation}
    \mathcal{C}(V', E')
    \;=\;
    \sum_{\ell \in V_{\text{logit}}} \mathrm{Rel}_{V'}(\ell)
    \;\in\; [0, 1]
    \label{eq:Ccal}
    \end{equation}
    is the fraction of source mass that still reaches a logit after
    pruning. We report $L_{\text{prune}} = 1 - \mathcal{C}$ alongside $L$
    of Eq.~\eqref{eq:Lsum} but do not include it in the simplex sum:
    $L_{\text{prune}}$ depends only on $(V', E')$ and not on $\varphi$, so
    the two-stage pipeline lets it govern Stage~1 cleanly while
    Eq.~\eqref{eq:Lsum} governs Stage~2 on a fixed pruned graph.